\subsection{Miscellaneous}
\label{appx:analyses:miscellaneous}
In this section, we include additional minor analyses around agent behavior and their generations.

\textbf{Agents are better at localizing files than BM25.}
The interactive setting also enables agents to identify the correct file(s) to edit more often compared to the RAG baselines in \citet{jimenez2024swebench}.
To measure this, we calculate the F1 score between the set of [edited, removed] files by the agent's prediction versus the gold patch.
SWE-agent w/ GPT-4 Turbo achieves an F1 score of \num{59.05}\%, while BM25 w/ Claude 3 Opus produces an F1 score of just \num{45.47}\%.

\textbf{Most resolved task instances are intentionally submitted.} There are four ways a task episode ends. 
\begin{itemize}
    \item ``Submit" refers to a task episode that ends when the agent generates the \texttt{submit} command.
    \item ``Exit Cost (Submit)" refers to the scenario where the episode ends because the cost limit was hit, and the changes so far are gathered and submitted as an edit.
    \item ``Exit Cost (No Submit)" refers to when the cost limit was hit and no \texttt{edit}'s were made, so there was nothing to submit. In this scenario, the instance is guaranteed to be unresolved.
    \item ``Early Exit" refers to when the task episode terminates because an agent issued too many malformed responses in a row. Any changes are submitted as an edit.
\end{itemize}

Table~\ref{tab:appx:exit_conditions} shows the counts for the number of trajectories that ended on these four different outcomes, categorized across the agent, SWE-bench split, and whether or not that task instance was resolved.
For SWE-agent with GPT-4 Turbo, the majority of ``All" task instances are \texttt{submit}.
For the trajectories corresponding to``All" task instances by SWE-agent with Claude 3 Opus, slightly less than \num{50}\% of task instances are submitted, while the slight majority are auto-submitted when the cost limit is hit.

\begin{table}[ht]
    \caption{
    This table showcases the counts for the four ways (``Submit", ``Exit Cost (Submit)", ``Exit Cost (No Submit)", ``Early Exit") a task episode could conclude.
    }
    \label{tab:appx:exit_conditions}
    \vspace{0.25em}
    \centering
    \begin{tabular}{lll|cccc}
         \toprule
            & & & Submit & Exit Cost & Exit Cost & Early Exit \\
            Model & Split & Outcome & & (Submit) & (No Submit) & \\
         \midrule
         SWE-agent                  & Full & Resolved & 266  & 20  & 0  & 0 \\
         \smalltab w/ GPT-4 Turbo   &      & All      & 1589 & 630 & 48 & 1 \\
                                    & Lite & Resolved & 50   & 4   & 0  & 0 \\
                                    &      & All      & 203  & 95  & 2  & 0 \\
         \midrule
         SWE-agent                  & Full & Resolved & 206  & 35   & 0  & 0 \\
         \smalltab w/ Claude 3 Opus &      & All      & 882  & 1048 & 73 & 1 \\
                                    & Lite & Resolved & 32   & 3    & 0  & 0 \\
                                    &      & All      & 133  & 156  & 11 & 0 \\
         \bottomrule
    \end{tabular}
\end{table}

However, these trends do not hold for ``Resolved" task instances.
For SWE-agent with both models, the large majority of these task instances are \texttt{submit}.
Reiterating the conclusion in Section~\ref{sec:analysis:agent-behavior} and prior visualizations in Section~\ref{appx:analyses:trajectory_analysis}, we see here again that resolved task instances often imply that the agent is able to produce and verify an edit within the allotted number of turns.
The SWE-agent ACI is also effective at eliciting well-formed thoughts and actions from agents.
Across all runs, there are only two ``Early Exit" occurrences, where the episode terminated because the agent generated too many malformed responses in a row.

Finally, Table~\ref{tab:appx:exit_conditions} also upholds an expected trend.
Task instances that finish with a \texttt{submit} action are more likely to be resolved than those that are cutoff by cost.
For instance, for SWE-agent with GPT-4 Turbo on full SWE-bench, \num{14.3}\% of task instances that end with a \texttt{submit} are resolved, which is much higher than \num{3.1}\% for those finishing on \texttt{exit\_cost}.
